\documentclass[11pt]{article}
\usepackage{graphicx} % Required for inserting images
\usepackage{xcolor}
\usepackage{fancybox}
\usepackage{amsmath}
\usepackage{natbib, graphicx, color, tabularx, mathtools, bbold, soul, bbm, enumerate, eurosym, multirow, booktabs}

\usepackage{natbib, graphicx, color, tabularx, mathtools, bbold, soul, bbm, enumerate, eurosym, multirow, booktabs}
\usepackage{amsfonts}



\usepackage[margin=0.8in]{geometry}
\usepackage{amsmath}
\usepackage{amssymb}
\usepackage{xcolor}
\usepackage{array}
\usepackage{enumitem}

\begin{document}
\begin{center}
{\Large The Hong Kong University of Science and Technology}\\
{\Large IEDA 1250, Summer  2026}\\
{\Large Mid-term Exam Solution}
\end{center}

\vspace{0.5cm}

\noindent
Student ID \underline{\hspace{3cm}}
\hfill
Name: \underline{\hspace{6cm}}

\vspace{0.8cm}

\begin{center}
\setlength{\fboxsep}{8pt}
\doublebox{
\begin{minipage}{0.80\textwidth}
\begin{center}
\textbf{Please read the following instructions carefully.}
\end{center}

\begin{itemize}
    \item You have \textbf{90 minutes} to complete this exam. 
    This question booklet contains 6 questions (including a bonus question) 
    for the total of 120 points. Check to see if any pages are missing.

    \item Read the instructions for individual questions carefully before answering 
    the questions.
\end{itemize}
\end{minipage}
}
\end{center}

\vspace{1cm}

\begin{center}
\renewcommand{\arraystretch}{2.0}
\begin{tabular}{|c|c|c|}
\hline
Question & Points & Score \\
\hline
1 & 10 & \\
\hline
2 & 20 & \\
\hline
3 & 20 & \\
\hline
4 & 20 & \\
\hline
5 & 20 & \\
\hline
6 & 30 & \\
\hline
Total & 120 & \\
\hline
\end{tabular}
\end{center}

\newpage
\begin{enumerate}

% =========================
% Question 1
% =========================

\item \textbf{(10 points)} Solve the following linear program (LP) graphically. Clearly identify the feasible region, the optimal solution, and the optimal objective value.
\[
\begin{aligned}
\max_{x_1,x_2}\quad & x_1+\frac{1}{2}x_2\\
\text{s.t.}\quad 
& x_1+2x_2\leq 6,\\
& x_1+2x_2\geq 0,\\
& 2x_1-x_2\leq 4,\\
& 2x_1-x_2\geq -2.
\end{aligned}
\]

{\color{red}
\textbf{Solution:}
The feasible region is a diamond with vertices
\[
\left(-\frac{4}{5},\frac{2}{5}\right),
\quad
\left(\frac{2}{5},\frac{14}{5}\right),
\quad
\left(\frac{14}{5},\frac{8}{5}\right),
\quad
\left(\frac{8}{5},-\frac{4}{5}\right).
\]

Now evaluate the objective function
\[
z=x_1+\frac{1}{2}x_2
\]
at the four vertices:
\[
\begin{array}{c|c}
\text{Vertex} & z=x_1+\frac{1}{2}x_2\\
\hline
\left(-\frac{4}{5},\frac{2}{5}\right) & -\frac{3}{5}\\[0.2cm]
\left(\frac{8}{5},-\frac{4}{5}\right) & \frac{7}{5}\\[0.2cm]
\left(\frac{2}{5},\frac{14}{5}\right) & \frac{9}{5}\\[0.2cm]
\left(\frac{14}{5},\frac{8}{5}\right) & \frac{18}{5}
\end{array}
\]
Therefore, the optimal solution is
\[
\boxed{
x_1^*=\frac{14}{5},
\qquad
x_2^*=\frac{8}{5}
}
\]
and the optimal objective value is
\[
\boxed{
z^*=\frac{18}{5}.
}
\]
}

\vspace{1em}

% =========================
% Question 2
% =========================

\item \textbf{(20 points)} Consider the following primal LP:
\[
\begin{aligned}
\max \quad & 2x_1+3x_2+x_3+x_4+x_5\\
\text{s.t.}\quad
& 2x_1+x_2+x_3+x_4\leq 3,\\
& x_1-x_2+2x_3+x_5=7,\\
& x_1+2x_2-x_3+3x_4+x_5\leq 6,\\
& x_2+x_3-x_4+2x_5=7,\\
& x_1\geq 0,\quad x_2\geq 0,
\end{aligned}
\]
where $x_3,x_4,x_5$ are unrestricted in sign. Suppose that $y_1,y_2,y_3,$ and $y_4$ are the dual variables corresponding to the four constraints of the primal LP, respectively, and that the optimal dual solution is
\[
(y_1^*,y_2^*,y_3^*,y_4^*)=(2,-1,0,1).
\]
Use the complementary slackness conditions to find an optimal primal solution.

{\color{red}
\textbf{Solution:}
Given the primal LP and the optimal dual solution
\[
(y_1^*,y_2^*,y_3^*,y_4^*)=(2,-1,0,1),
\]
we use complementary slackness to find an optimal primal solution.

The dual LP is
\[
\begin{aligned}
\min \quad & 3y_1+7y_2+6y_3+7y_4\\
\text{s.t.}\quad
& 2y_1+y_2+y_3 \geq 2,\\
& y_1-y_2+2y_3+y_4 \geq 3,\\
& y_1+2y_2-y_3+y_4 = 1,\\
& y_1+3y_3-y_4 = 1,\\
& y_2+y_3+2y_4 = 1,\\
& y_1\geq 0,\quad y_3\geq 0,\quad y_2,y_4 \text{ free.}
\end{aligned}
\]

Step 1: Apply complementary slackness.

For primal constraints:
\[
y_1^*\left(3-(2x_1+x_2+x_3+x_4)\right)=0.
\]
Since $y_1^*=2\neq 0$, the first primal constraint must be tight:
\[
2x_1+x_2+x_3+x_4=3.
\]
Also,
\[
y_3^*\left(6-(x_1+2x_2-x_3+3x_4+x_5)\right)=0.
\]
Since $y_3^*=0$, this gives no additional equality condition.

For dual constraints corresponding to $x_1$ and $x_2$:
\[
x_1\left(2y_1^*+y_2^*+y_3^*-2\right)=0.
\]
Substituting the dual solution gives
\[
2(2)+(-1)+0-2=1\neq 0.
\]
Therefore,
\[
x_1=0.
\]
Similarly,
\[
x_2\left(y_1^*-y_2^*+2y_3^*+y_4^*-3\right)=0.
\]
Substituting the dual solution gives
\[
2-(-1)+2(0)+1-3=1\neq 0.
\]
Therefore,
\[
x_2=0.
\]

Step 2: Substitute $x_1=0$ and $x_2=0$ into the primal constraints.

The tight first primal constraint gives
\[
x_3+x_4=3.
\]
The second primal constraint gives
\[
2x_3+x_5=7.
\]
The fourth primal constraint gives
\[
x_3-x_4+2x_5=7.
\]
Therefore, we solve
\[
\begin{cases}
x_3+x_4=3,\\
2x_3+x_5=7,\\
x_3-x_4+2x_5=7.
\end{cases}
\]
Solving this system, we obtain
\[
x_3=2,\qquad x_4=1,\qquad x_5=3.
\]
Hence, an optimal primal solution is
\[
\boxed{
x_1^*=0,\quad
x_2^*=0,\quad
x_3^*=2,\quad
x_4^*=1,\quad
x_5^*=3.
}
\]
The optimal primal objective value is
\[
2x_1^*+3x_2^*+x_3^*+x_4^*+x_5^*
=
0+0+2+1+3
=
6.
\]
Thus,
\[
\boxed{z^*=6.}
\]
}

\vspace{1em}

% =========================
% Question 3
% =========================

\item \textbf{(20 points)} A relief agency needs to ship medicine boxes from two warehouses to three clinics. Warehouse 1 has 20 boxes available, and Warehouse 2 has 30 boxes available. Clinic 1 needs 10 boxes, Clinic 2 needs 25 boxes, and Clinic 3 needs 15 boxes. Since total supply equals total demand, all boxes must be shipped. The shipping cost per box is shown below.
\[
\begin{array}{c|ccc}
& \text{Clinic 1} & \text{Clinic 2} & \text{Clinic 3}\\
\hline
\text{Warehouse 1} & 4 & 6 & 8\\
\text{Warehouse 2} & 5 & 4 & 3
\end{array}
\]
Let $x_{ij}$ denote the number of boxes shipped from Warehouse $i$ to Clinic $j$.

\begin{enumerate}
    \item[(a)] Formulate a linear program that minimizes the total shipping cost while satisfying all clinic demands.

    \item[(b)] Write down the dual linear program. Let $u_1,u_2$ be the dual variables associated with the two warehouse constraints, and let $v_1,v_2,v_3$ be the dual variables associated with the three clinic constraints.

    \item[(c)] Suppose the optimal shipment plan is
    \[
    x_{11}^*=10,\quad x_{12}^*=10,\quad x_{13}^*=0,\quad
    x_{21}^*=0,\quad x_{22}^*=15,\quad x_{23}^*=15.
    \]
    Also suppose an optimal dual solution is
    \[
    u_1^*=2,\quad u_2^*=0,\quad v_1^*=2,\quad v_2^*=4,\quad v_3^*=3.
    \]
    Explain the meaning of the dual variables intuitively. What is the economic interpretation of the dual LP and its optimal solution?
\end{enumerate}

{\color{red}
\textbf{Solution:}
Let $x_{ij}$ denote the number of boxes shipped from Warehouse $i$ to Clinic $j$.

\textbf{(a) Primal LP formulation.}

The objective is to minimize the total shipping cost:
\[
\min \quad
4x_{11}+6x_{12}+8x_{13}
+5x_{21}+4x_{22}+3x_{23}.
\]
Subject to the warehouse supply constraints:
\[
x_{11}+x_{12}+x_{13}=20,
\]
\[
x_{21}+x_{22}+x_{23}=30,
\]
and the clinic demand constraints:
\[
x_{11}+x_{21}=10,
\]
\[
x_{12}+x_{22}=25,
\]
\[
x_{13}+x_{23}=15.
\]
Also,
\[
x_{ij}\geq 0,
\qquad
i=1,2,\quad j=1,2,3.
\]

\textbf{(b) Dual LP formulation.}

Let $u_1,u_2$ be the dual variables associated with the two warehouse constraints, and let $v_1,v_2,v_3$ be the dual variables associated with the three clinic constraints. Since all primal constraints are equality constraints, all dual variables are unrestricted in sign.

The dual LP is
\[
\max \quad 20u_1+30u_2+10v_1+25v_2+15v_3
\]
subject to
\[
u_1+v_1\leq 4,
\]
\[
u_1+v_2\leq 6,
\]
\[
u_1+v_3\leq 8,
\]
\[
u_2+v_1\leq 5,
\]
\[
u_2+v_2\leq 4,
\]
\[
u_2+v_3\leq 3,
\]
where
\[
u_1,u_2,v_1,v_2,v_3 \text{ are free.}
\]

\textbf{(c) Economic interpretation of the dual variables.}

The dual variables can be interpreted as shadow prices. The variables $u_1$ and $u_2$ represent the marginal values associated with the supplies available at Warehouses 1 and 2, respectively. The variables $v_1,v_2,v_3$ represent the marginal values associated with satisfying demands at Clinics 1, 2, and 3, respectively.

The dual constraint
\[
u_i+v_j\leq c_{ij}
\]
means that the combined value of one box supplied from Warehouse $i$ and delivered to Clinic $j$ cannot exceed the actual shipping cost $c_{ij}$. Therefore, the dual LP tries to assign values to warehouse supplies and clinic demands so that these values are consistent with all shipping costs.

Using the given optimal dual solution
\[
u_1^*=2,\qquad u_2^*=0,\qquad v_1^*=2,\qquad v_2^*=4,\qquad v_3^*=3,
\]
the dual objective value is
\[
20(2)+30(0)+10(2)+25(4)+15(3)
=
40+0+20+100+45
=
205.
\]
The cost of the given primal shipment plan is
\[
4(10)+6(10)+8(0)+5(0)+4(15)+3(15)
=
40+60+0+0+60+45
=
205.
\]
Since the primal and dual objective values are equal, the given shipment plan and the given dual solution are both optimal.

Economically, the optimal dual solution provides a system of prices that supports the optimal transportation plan. Routes with positive shipments satisfy
\[
u_i^*+v_j^*=c_{ij},
\]
while routes with zero shipments satisfy
\[
u_i^*+v_j^*\leq c_{ij}.
\]
Thus, only routes whose shipping cost is justified by the corresponding warehouse and clinic shadow prices are used in the optimal plan.
}

\vspace{1em}

% =========================
% Question 4
% =========================

\item \textbf{(20 points)} A manufacturer produces a single product over the next 6 months. The demand in each month is known and must be satisfied on time. Shortages are not allowed. The manufacturer may produce in a month only if it sets up the production line in that month. If the production line is set up in month $t$, where $t\in\{1,2,\ldots,6\}$, the firm pays a fixed setup cost. In addition, it pays a variable production cost for each unit produced and an inventory holding cost for each unit carried from one month to the next.

The monthly data are shown below.
\[
\begin{array}{c|cccccc}
\text{Month} & 1 & 2 & 3 & 4 & 5 & 6\\
\hline
\text{Demand} & 40 & 60 & 35 & 70 & 55 & 45\\
\text{Production capacity} & 80 & 80 & 80 & 80 & 80 & 80\\
\text{Variable production cost} & 12 & 11 & 13 & 10 & 12 & 14\\
\text{Setup cost} & 300 & 300 & 300 & 300 & 300 & 300\\
\text{Holding cost} & 2 & 2 & 2 & 2 & 2 & -
\end{array}
\]
Assume that the initial inventory is zero and that the company wants no inventory left over at the end of Month 6. Formulate this problem as a mixed-integer linear program that minimizes the total production, setup, and inventory holding cost. Clearly define all decision variables.

{\color{red}
\textbf{Solution:}
Let $x_t$ be the number of units produced in month $t$, $I_t$ be the inventory at the end of month $t$, and
\[
y_t=
\begin{cases}
1, & \text{if the production line is set up in month }t,\\
0, & \text{otherwise,}
\end{cases}
\qquad t=1,\ldots,6.
\]

The mixed-integer linear program is
\[
\begin{aligned}
\min \quad
&12x_1+11x_2+13x_3+10x_4+12x_5+14x_6\\
&+300(y_1+y_2+y_3+y_4+y_5+y_6)\\
&+2(I_1+I_2+I_3+I_4+I_5)
\end{aligned}
\]
subject to
\[
x_1=40+I_1,
\]
\[
I_1+x_2=60+I_2,\qquad
I_2+x_3=35+I_3,
\]
\[
I_3+x_4=70+I_4,\qquad
I_4+x_5=55+I_5,
\]
\[
I_5+x_6=45+I_6,
\]
\[
I_6=0,
\]
\[
0\leq x_t\leq 80y_t,\qquad t=1,\ldots,6,
\]
\[
I_t\geq 0,\qquad t=1,\ldots,6,
\]
and
\[
y_t\in\{0,1\},\qquad t=1,\ldots,6.
\]
The constraints $x_t\leq 80y_t$ ensure that production in month $t$ is possible only if the setup decision $y_t=1$.
}

\vspace{1em}

% =========================
% Question 5
% =========================

\item \textbf{(20 points)} Consider a two-player zero-sum game where Player A's payoffs are shown below. Player A chooses rows $(A1$ to $A5)$, and Player B chooses columns $(B1$ to $B5)$.
\[
\begin{array}{c|ccccc}
\text{Player A}\backslash \text{Player B} & B1 & B2 & B3 & B4 & B5\\
\hline
A1 & 4 & 1 & 5 & 2 & 6\\
A2 & 2 & 3 & 4 & 4 & 5\\
A3 & 1 & 0 & 2 & 1 & 3\\
A4 & 1 & 2 & 3 & 3 & 4\\
A5 & 0 & 0 & 1 & 1 & 2
\end{array}
\]
Answer the following questions:
\begin{enumerate}
    \item[(a)] Eliminate all dominated strategies for Player A and Player B.
    \item[(b)] Does the reduced game have a pure strategy Nash equilibrium? If yes, find the equilibrium; if no, explain why every payoff in the reduced game is unstable.
\end{enumerate}

{\color{red}
\textbf{Solution:}
Since this is a zero-sum game, Player A maximizes the payoff and Player B minimizes Player A's payoff.

\textbf{(a) Elimination of dominated strategies.}
For Player A, Row $A3$ is dominated by Row $A1$, Row $A4$ is dominated by Row $A2$, and Row $A5$ is dominated by Row $A1$. Therefore, Rows $A3,A4,A5$ can be eliminated.

For Player B, Column $B5$ is dominated by Column $B3$, Column $B3$ is dominated by Column $B2$, and Column $B4$ is dominated by Column $B2$. Therefore, Columns $B3,B4,B5$ can be eliminated.

The reduced game is
\[
\begin{array}{c|cc}
 & B1 & B2\\
\hline
A1 & 4 & 1\\
A2 & 2 & 3
\end{array}
\]

\textbf{(b) Pure strategy Nash equilibrium.}
The row minima are
\[
\min(A1)=1,\qquad \min(A2)=2,
\]
so $\max\min=2$. The column maxima are
\[
\max(B1)=4,\qquad \max(B2)=3,
\]
so $\min\max=3$. Since
\[
\max\min=2\neq 3=\min\max,
\]
there is no pure strategy Nash equilibrium.

Equivalently, every payoff in the reduced game is unstable:
\begin{itemize}
    \item At $(A1,B1)$, Player B can switch to $B2$ and reduce A's payoff from $4$ to $1$.
    \item At $(A1,B2)$, Player A can switch to $A2$ and increase the payoff from $1$ to $3$.
    \item At $(A2,B1)$, Player A can switch to $A1$ and increase the payoff from $2$ to $4$.
    \item At $(A2,B2)$, Player B can switch to $B1$ and reduce A's payoff from $3$ to $2$.
\end{itemize}
Therefore, the reduced game has no pure strategy Nash equilibrium.
}

\vspace{1em}

% =========================
% Question 6
% =========================

\item \textbf{(30 points)} A firm has options to invest in 10 different projects. The total investment cannot exceed the available budget $q$. The profits and costs of these projects are given by $(p_1,p_2,\ldots,p_{10})$ and $(c_1,c_2,\ldots,c_{10})$, respectively. The firm wants to maximize its total profit. In addition to the budget constraint, the managers need to consider the following requirements:
\begin{enumerate}
    \item[(a)] Projects 2 and 5 cannot be chosen together.
    \item[(b)] Exactly 4 projects are to be selected.
    \item[(c)] At least one project must be selected from the project set $\{1,3,4\}$.
    \item[(d)] If Project 6 is selected, then Project 7 must also be selected.
    \item[(e)] If Project 8 is selected, then at least one of Projects 3 and 4 must also be selected.
    \item[(f)] Either both Projects 1 and 10 are selected, or neither of them is selected.
    \item[(g)] If both Projects 3 and 4 are selected, then Project 9 must also be selected.
    \item[(h)] Project 5 cannot be selected unless at least two projects from the set $\{6,7,8\}$ are selected.
    \item[(i)] At least one of the following two conditions must hold: at least two projects from the set $\{1,2,3,4\}$ are selected, or at most one project from the set $\{5,6,7,8\}$ is selected.
    \item[(j)] If Project 10 is selected, then at least one of the following two conditions must hold:
    \begin{itemize}
        \item Both Projects 2 and 9 are selected.
        \item Neither Project 3 nor Project 4 is selected.
    \end{itemize}
\end{enumerate}
Define integer decision variables and formulate logical constraints for (a)--(j).

{\color{red}
\textbf{Solution:}
Let
\[
x_i=
\begin{cases}
1, & \text{if Project }i\text{ is selected,}\\
0, & \text{otherwise,}
\end{cases}
\qquad i=1,\ldots,10.
\]
The objective is
\[
\max \quad \sum_{i=1}^{10}p_ix_i.
\]
The budget constraint is
\[
\sum_{i=1}^{10}c_ix_i\leq q,
\]
with
\[
x_i\in\{0,1\},\qquad i=1,\ldots,10.
\]

The logical constraints are
\[
x_2+x_5\leq 1,
\]
\[
\sum_{i=1}^{10}x_i=4,
\]
\[
x_1+x_3+x_4\geq 1,
\]
\[
x_6\leq x_7,
\]
\[
x_8\leq x_3+x_4,
\]
\[
x_1=x_{10},
\]
\[
x_9\geq x_3+x_4-1,
\]
\[
x_6+x_7+x_8\geq 2x_5.
\]

For part (i), introduce a binary variable $z_i\in\{0,1\}$ and impose
\[
x_1+x_2+x_3+x_4\geq 2z_i,
\]
\[
x_5+x_6+x_7+x_8\leq 1+3z_i.
\]
If $z_i=1$, the first condition is enforced. If $z_i=0$, the second condition is enforced.

For part (j), introduce a binary variable $z_j\in\{0,1\}$ and impose
\[
x_2\geq x_{10}+z_j-1,
\]
\[
x_9\geq x_{10}+z_j-1,
\]
\[
x_3\leq 1-x_{10}+z_j,
\]
\[
x_4\leq 1-x_{10}+z_j.
\]
If $x_{10}=1$ and $z_j=1$, then $x_2=x_9=1$. If $x_{10}=1$ and $z_j=0$, then $x_3=x_4=0$. If $x_{10}=0$, these constraints impose no additional restriction.
}

\end{enumerate}

\end{document}
